% Checkpoint 1 report. The numbers, tables and figure come from report_numbers.tex,
% which white_advantage writes when run with --report (see scripts/README.md).
% Build from the repo root:  pdflatex -jobname=checkpointOne report.tex
\documentclass[11pt,a4paper]{article}
\usepackage[a4paper,margin=1.8cm,top=1.5cm,bottom=1.6cm]{geometry}
\usepackage[T1]{fontenc}
\usepackage[sc]{mathpazo}
\linespread{1.03}
\usepackage{courier}
\usepackage[scaled=0.9]{helvet}
\usepackage{microtype}
\usepackage{xcolor}
\usepackage{graphicx}
\usepackage{booktabs}
\usepackage{tabularx}
\usepackage{enumitem}
\usepackage{amsmath}
\usepackage[hidelinks]{hyperref}
\usepackage{titlesec}
\usepackage{pgfplots}
\pgfplotsset{compat=1.16}

\titleformat{\section}{\large\sffamily\bfseries}{\thesection.}{0.5em}{}
\titleformat{\subsection}{\normalsize\sffamily\bfseries}{\thesubsection}{0.5em}{}
\titlespacing*{\section}{0pt}{8pt}{3pt}
\titlespacing*{\subsection}{0pt}{5pt}{2pt}
\setlist{itemsep=1pt,topsep=2pt,leftmargin=1.2em}
\setlength{\parindent}{0pt}
\setlength{\parskip}{3pt}
\pagestyle{empty}

% red placeholder for anything not filled in yet
\newcommand{\TBD}[1]{\textcolor{red}{\textbf{[#1]}}}

% placeholder: white_advantage --report report_numbers.tex replaces this file with the real numbers
\newcommand{\MonthList}{\TBD{MonthList}}
\newcommand{\GamesRead}{\TBD{GamesRead}}
\newcommand{\GamesUsed}{\TBD{GamesUsed}}
\newcommand{\PctUsed}{\TBD{PctUsed}}
\newcommand{\SkipNoResult}{\TBD{SkipNoResult}}
\newcommand{\SkipAbandoned}{\TBD{SkipAbandoned}}
\newcommand{\SkipNoRating}{\TBD{SkipNoRating}}
\newcommand{\SkipGap}{\TBD{SkipGap}}
\newcommand{\ScoreAll}{\TBD{ScoreAll}}
\newcommand{\ExpectedAll}{\TBD{ExpectedAll}}
\newcommand{\HAll}{\TBD{HAll}}
\newcommand{\HAllCI}{\TBD{HAllCI}}
\newcommand{\MainSpeed}{\TBD{MainSpeed}}
\newcommand{\MainBands}{\TBD{MainBands}}
\newcommand{\HMin}{\TBD{HMin}}
\newcommand{\HMinBand}{\TBD{HMinBand}}
\newcommand{\HMax}{\TBD{HMax}}
\newcommand{\HMaxBand}{\TBD{HMaxBand}}
\newcommand{\FileRows}{\multicolumn{4}{c}{\TBD{filled in by white\_advantage}} \\}
\newcommand{\SpeedRows}{\multicolumn{5}{c}{\TBD{filled in by white\_advantage}} \\}
\newcommand{\Figure}{\fbox{\parbox[c][3.5cm][c]{0.95\linewidth}{\centering\TBD{figure: h by rating bracket, drawn by white\_advantage}}}}


% fill these in by hand
\newcommand{\Mentor}{Namish Bansal}
\newcommand{\Repo}{\TBD{GitHub repository link (group repo from Classroom50)}}

\begin{document}

\begin{center}
{\Large\bfseries The Perks of Being a First Mover}\\[2pt]
{\large Does White's Advantage in Chess Grow or Shrink with Rating?}\\[2pt]
CS215 Open Project \quad$\cdot$\quad Checkpoint 1
\end{center}

\section{Meta Information}
\begin{tabularx}{\linewidth}{@{}lX@{}}
\textbf{Project title} & The Perks of Being a First Mover \\
\textbf{Group number} & 02 \\
\textbf{Mentor} & \Mentor \\
\textbf{Members} & Sanidhya (25b0944), Daksh Maahor (25b0974), Sahil (25b0991), VadaPavJai (25b1030) \\
\textbf{Repository} & \Repo \\
\textbf{Tracking sheet} & \scriptsize\url{https://docs.google.com/spreadsheets/d/1pvTllHR4a1Sl7GHtjKnvLeb-zmmHaM3BjltyaDPX9-Q} \\
\end{tabularx}

\section{Progress So Far}

\subsection*{Dataset and descriptive statistics}
\textbf{Source:} Lichess Open Database (\url{https://database.lichess.org}), standard rated games for
\MonthList, released under the Creative Commons CC0 licence, so it is free to download, use and
redistribute. Each month is one zstd-compressed PGN file of tens of GB (well over 100 GB once
decompressed). Every game is one record with about 17 header tags (players, ratings, result, date,
opening, time control, termination) followed by the move list. We use 6 tags: \texttt{Result},
\texttt{WhiteElo}, \texttt{BlackElo}, \texttt{Event} (gives the speed), \texttt{TimeControl} and
\texttt{Termination}.

\begin{center}
\small
\begin{tabular}{@{}lrrr@{}}
\toprule
Month & Games read & Games kept & Kept \\
\midrule
\FileRows
\midrule
Total & \GamesRead & \GamesUsed & \PctUsed \\
\bottomrule
\end{tabular}
\hspace{1.2em}
\begin{tabular}{@{}lrccc@{}}
\toprule
Speed & Games & White score & Expected & $h$ (pts, 95\% CI) \\
\midrule
\SpeedRows
\bottomrule
\end{tabular}
\end{center}

\emph{White score} is White's average score (win 1, draw $\tfrac12$, loss 0); \emph{Expected} is the
score the two ratings alone predict. Over all kept games White scores \ScoreAll{} against an expected
\ExpectedAll, which corresponds to $h = \HAll \pm \HAllCI$ rating points. Within \MainSpeed{}, $h$
ranges from \HMin{} in bracket \HMinBand{} to \HMax{} in \HMaxBand{} across \MainBands{} brackets
(figure below; brackets with fewer than 1{,}000 games are hidden).

\begin{center}
\Figure
\end{center}

\subsection*{Data cleaning steps completed}
\begin{itemize}
\item \textbf{Parser:} a C++ program (\texttt{scripts/white\_advantage.cpp}) streams each
  compressed file through the zstd library in 1\,MB chunks, reads only header lines and skips the
  moves. Memory stays at a few MB, and the two months are processed in parallel, one thread each.
 
\item \textbf{Removed:} games with no result (\SkipNoResult), abandoned games (\SkipAbandoned),
  games with a missing or ``?'' rating (\SkipNoRating), and games where the rating gap is over
  100 points (\SkipGap).
\item \textbf{Grouped:} by speed (UltraBullet to Correspondence, from the \texttt{Event} tag) and by
  100-point rating bracket $[a, a{+}100)$ of the \emph{average} of both ratings, so every game falls
  in exactly one bracket.
\item \textbf{Aggregated:} instead of exporting one CSV row per game, we store game counts and
  White's total score for every (speed, bracket, rating gap) cell. This is all the model needs, and
  it fits in a few MB.
\item \textbf{Checked:} on synthetic games with a known built-in advantage of $+35$ points, the
  pipeline recovers $+34.8 \pm 1.7$.
\end{itemize}

\subsection*{Division and revision of responsibilities}
\begin{tabularx}{\linewidth}{@{}lX@{}}
\toprule
Member & Responsibility (revised since the proposal: work split equally) \\
\midrule
Sanidhya & Statistical modelling: fitting $h$ and its confidence intervals; next, the trend test across brackets. \\
Daksh Maahor & Data cleaning: filters, speed and rating-bracket grouping; next, separate win/draw rates per bracket. \\
Sahil & Data engineering: C++ streaming parser for the \texttt{.pgn.zst} files; next, robustness runs across time controls. \\
VadaPavJai & Visualisation and report: figures and this report; next, the Plotly dashboard. \\
\bottomrule
\end{tabularx}

\subsection*{Refined problem statement}
The proposal compared raw win rates across Elo brackets. Raw win rates are confounded by the rating
gap: a 1540 who beats a 1570 with White is stronger evidence of a first-move edge than a win between
two 1570s. We therefore model each game with the Elo expected-score curve plus a White bonus $h$:
\[
E = \frac{1}{1 + 10^{-(R_\text{White} + h - R_\text{Black})/400}},
\]
and estimate $h$, White's advantage in rating points, by maximum likelihood separately for each speed
and each 100-point bracket, with a 95\% confidence interval from the likelihood curvature. Draws count
as half a point. The question becomes: \emph{does $h$ rise to a peak at intermediate ratings and
fall at 2200+?} We also widened the rating-gap filter from 50 to 100 points, since the model now
corrects for the gap itself and the wider filter keeps far more games.

For Checkpoint 2: a formal test for a trend in $h$ across brackets, separate win and draw rates per
bracket (draws rise with rating), robustness across time controls, and a Plotly dashboard.

\subsection*{Tracking spreadsheet}
Updated with Checkpoint 1 rows (meta information, dataset, cleaning, revised roles, refined problem
statement); the mentor has edit access. Link in Section 1.

\subsection*{Blockers and challenges}
\begin{itemize}
\item \textbf{Raw data can't go on GitHub.} The guidelines ask us to push data files, but each
  monthly file is tens of GB, far above GitHub's limits (100 MB per file, 2 GB with Git LFS on the
  free plan). We are pushing the aggregated output (\texttt{output/summary.txt},
  \texttt{output/output.csv}) and the exact download links instead. Is this acceptable?
\item The two concerns from our proposal are resolved: streaming the files avoids the memory limit,
  and draws are handled as half a point inside the model.
\end{itemize}

\end{document}
